problem:  
Let \((M,\mathcal{D},g)\) be a **compact, equiregular step‑2 sub‑Riemannian manifold** equipped with its Popp measure and associated symmetric sub‑Laplacian \(\Delta\). Denote by \(p_t(x,y)\) the heat kernel of \(\Delta\) and by \(\mathfrak{g}_x\) the nilpotent (Carnot) tangent Lie algebra at \(x\), determined via nilpotentization of \((\mathcal{D},[\mathcal{D},\mathcal{D}])\).  

It is known that, as \(t\to0^+\),
\[
p_t(x,x)\sim \frac{1}{t^{Q/2}}\big(1+a_1(x)t+a_2(x)t^2+\cdots\big),
\]
where \(Q\) is the homogeneous dimension and the coefficients \(a_j(x)\) depend polynomially on the curvature‑type tensors and their derivatives arising from the sub‑Riemannian connection associated with \((\mathcal{D},g)\).  

**Open problem:**  
Assume that, for all \(x\in M\), the function
\[
a_j(x)\equiv a_j^{\mathbb{G}}\quad\text{for all } j\in\mathbb{N},
\]
where each \(a_j^{\mathbb{G}}\) equals the corresponding coefficient for some fixed *nonabelian* step‑2 Carnot group \(\mathbb{G}\) (so the entire on‑diagonal heat kernel expansion of \((M,\Delta)\) matches that of the homogeneous model \(\mathbb{G}\)).  

Does it follow that, locally, each nilpotent tangent algebra \(\mathfrak{g}_x\) is **isomorphic** to the Lie algebra of \(\mathbb{G}\), and that \((M,\mathcal{D},g)\) is locally sub‑Riemannian *isometric* to an open subset of \(\mathbb{G}\)?  
Equivalently: do complete coincidence of the heat coefficients with those of a fixed Carnot group enforce local geometric homogeneity of the sub‑Riemannian structure?

Why is it a "good" problem:  
This formulation eliminates the previous logical issues and addresses a genuine rigidity question that lies precisely on the frontier of current knowledge. The small‑time asymptotic coefficients \(a_j(x)\) encode, via microlocal heat kernel parametrix constructions, all curvature‑type invariants arising from the brackets of \(\mathcal{D}\). It is known that in step‑2 structures these coefficients depend on finite jets of those invariants, yet it is not known whether equality of *all* coefficients with those of a fixed homogeneous model forces local homogeneity of the geometry.  

This is a profound inverse‑spectral rigidity question: it asks whether the entire analytic signature of the heat kernel uniquely encodes the local Carnot type. The problem unites **sub‑Riemannian microlocal analysis** (through the heat kernel asymptotics) and **differential geometry of nonholonomic distributions** (through curvature and nilpotentization) and targets one of the deepest open interfaces between these fields.  

A positive resolution would establish that the full infinite sequence of heat coefficients captures the complete local geometric structure, achieving a form of *infinitesimal spectral completeness* in the sub‑Riemannian setting. A negative resolution would demonstrate the existence of **isospectral but non‑isomorphic** sub‑Riemannian geometries, opening new directions in the study of geometric and spectral invariants beyond the Riemannian paradigm.  

Its statement is simple but its implications are vast: “Does perfect equality of all heat coefficients force local geometric equality?”—a paradigmatic example of a clear, deep, and cross‑disciplinary research problem.